\ifdefined\gkver\else\def\gkver{student}\fi
\documentclass[\gkver]{gaokaozhenti}
\newcommand{\ccnum}[1]{{\fontspec{Noto Sans CJK JP}#1}}
\begin{document}

\chapter{2019年全国理综Ⅰ}
%% paperType: 理综物理部分

\begin{enumerate}
\item
%% number: 14
%% paperName: 2019年全国理综Ⅰ第 14 题 6 分
%% typeId: 1
%% type: 单选题
%% chapter: 原子和原子核
%% point: 玻尔模型
%% method: 无
%% score: 6
%% degree: 950
%% duplicateId: 0
%% body:
氢原子能级示意图如图所示。光子能量在 $1.63\UeV\sim3.10\UeV$ 的光为可见光。要使处于基态（$n=1$）的氢原子被激发后可辐射出可见光光子，最少应给氢原子提供的能量为\xzanswer{A}

\onepicture[w=6.5cm, align=right]{\includesvg{figs/14}}

\fourchoices[answer=A]
{$12.09\UeV$}
{$10.20\UeV$}
{$1.89\UeV$}
{$1.51\UeV$}

%% answer: A
%% memo:
\memoanswer{
氢原子从 $n=3$ 能级向 $n=2$ 能级跃迁时释放的光子能量为 $\Delta E=E_3-E_2=1.89\UeV$，处于可见光光子能量范围，从 $n=4$ 能级向 $n=3$ 能级跃迁时释放的光子能量为 $\Delta E'=E_4-E_3=0.66\UeV$，不属于可见光范围，从 $n=2$ 能级向 $n=1$ 能级跃迁时释放的光子能量为 $\Delta E''=E_2-E_1=10.2\UeV$，不属于可见光范围，因此氢原子至少需从基态跃迁到 $n=3$ 能级才可辐射出可见光光子，则给氢原子提供的能量至少为 $\Delta E_1=E_3-E_1=12.09\UeV$，故 A 正确。
}

\item
%% number: 15
%% paperName: 2019年全国理综Ⅰ第 15 题 6 分
%% typeId: 1
%% type: 单选题
%% chapter: 电场
%% point: 带电粒子的平衡
%% method: 无
%% score: 6
%% degree: 650
%% duplicateId: 0
%% body:
如图，空间存在一方向水平向右的匀强电场，两个带电小球 P 和 Q 用相同的绝缘细绳悬挂在水平天花板下，两细绳都恰好与天花板垂直，则\xzanswer{D}

\onepicture[w=5.5cm, align=right]{\includesvg{figs/15}}

\fourchoices[answer=D]
{P 和 Q 都带正电荷}
{P 和 Q 都带负电荷}
{P 带正电荷，Q 带负电荷}
{P 带负电荷，Q 带正电荷}

%% answer: D
%% memo:
\memoanswer{
两绝缘细绳保持竖直状态，对两小球整体分析，整体受电场力的合力为零，电荷量为零，故两小球带等量异种电荷。隔离 P 进行研究，P 受库仑引力作用，则 P 受到的电场力由 Q 指向 P，可知 P 带负电荷，则 Q 带正电荷，故 D 正确。
}

\item
%% number: 16
%% paperName: 2019年全国理综Ⅰ第 16 题 6 分
%% typeId: 1
%% type: 单选题
%% chapter: 运动和力
%% point: 动量定理
%% method: 无
%% score: 6
%% degree: 750
%% duplicateId: 0
%% body:
最近，我国为“长征九号”研制的大推力新型火箭发动机联试成功，这标志着我国重型运载火箭的研发取得突破性进展。若某次实验中该发动机向后喷射的气体速度约为 $3\Ukms$，产生的推力约为 $4.8\times10^{6}\UN$，则它在 $1\Us$ 时间内喷射的气体质量约为\xzanswer{B}

\fourchoices[answer=B]
{$1.6\times10^{2}\Ukg$}
{$1.6\times10^{3}\Ukg$}
{$1.6\times10^{5}\Ukg$}
{$1.6\times10^{6}\Ukg$}

%% answer: B
%% memo:
\memoanswer{
设发动机单位时间内喷射气体的质量为 $\Delta m$，对喷出的气体进行研究有 $F\Delta t=\Delta mv$，代入数据解得 $\Delta m=1.6\times10^{3}\Ukg$，故 B 正确。
}

\item
%% number: 17
%% paperName: 2019年全国理综Ⅰ第 17 题 6 分
%% typeId: 1
%% type: 单选题
%% chapter: 磁场
%% point: 安培力
%% method: 无
%% score: 6
%% degree: 550
%% duplicateId: 0
%% body:
如图，等边三角形线框 LMN 由三根相同的导体棒连接而成，固定于匀强磁场中，线框平面与磁感应强度方向垂直，线框顶点 M、N 与直流电源两端相接，已知导体棒 MN 受到的安培力大小为 $F$，则线框 LMN 受到的安培力的大小为\xzanswer{B}

\onepicture[w=5.5cm, align=right]{\includesvg{figs/17}}

\fourchoices[answer=B]
{$2F$}
{$1.5F$}
{$0.5F$}
{$0$}

%% answer: B
%% memo:
\memoanswer{
设导体棒 MN 长度为 $l$，MN 棒中电流为 $I$，则其受到的安培力大小 $F=BIl$。MLN 的电阻是 MN 棒电阻的两倍，二者并联，两端电压相同，则流过 MLN 的电流为 $0.5I$，MLN 受到的安培力的合力为 $F_2=0.5BIl=0.5F$，MN 棒和 MLN 受到的安培力方向相同，故线框 LMN 受到的安培力大小为 $1.5F$，故 B 正确。
}

\item
%% number: 18
%% paperName: 2019年全国理综Ⅰ第 18 题 6 分
%% typeId: 1
%% type: 单选题
%% chapter: 直线运动
%% point: 竖直上下抛运动
%% method: 无
%% score: 6
%% degree: 550
%% duplicateId: 0
%% body:
如图，篮球架下的运动员原地垂直起跳扣篮，离地后重心上升的最大高度为 $H$。上升第一个 $\frac{H}{4}$ 所用的时间为 $t_1$，第四个 $\frac{H}{4}$ 所用的时间为 $t_2$。不计空气阻力，则 $\frac{t_2}{t_1}$ 满足\xzanswer{C}

\onepicture[w=4.5cm, align=right]{figs/18.png}

\fourchoices[answer=C]
{$1<\frac{t_2}{t_1}<2$}
{$2<\frac{t_2}{t_1}<3$}
{$3<\frac{t_2}{t_1}<4$}
{$4<\frac{t_2}{t_1}<5$}

%% answer: C
%% memo:
\memoanswer{
运动员的竖直上抛运动逆向等效为自由落体运动，即 $H=\frac{1}{2}gt^2$，全过程所用时间 $t=\sqrt{\frac{2H}{g}}$。下落第一个 $\frac{H}{4}$ 所用的时间为 $t_1$，满足 $\frac{H}{4}=\frac{1}{2}gt_1^2$，解得 $t_1'=\sqrt{\frac{H}{2g}}$；最后一个 $\frac{H}{4}$ 之前所用的时间为 $t_2$，满足 $\frac{3H}{4}=\frac{1}{2}gt_2^2$，解得 $t_2'=\sqrt{\frac{3H}{2g}}$。由此可知，$\frac{t_2}{t_1}=\frac{t_1'}{t-t_2}=\frac{\sqrt{\frac{H}{2g}}}{\sqrt{\frac{2H}{g}}-\sqrt{\frac{3H}{2g}}}$。故 C 正确。
}

\item
%% number: 19
%% paperName: 2019年全国理综Ⅰ第 19 题 6 分
%% typeId: 2
%% type: 多选题
%% chapter: 力物体的平衡
%% point: 动态平衡
%% method: 整体与隔离
%% score: 6
%% degree: 450
%% duplicateId: 0
%% body:
如图，一粗糙斜面固定在地面上，斜面顶端装有一光滑定滑轮。一细绳跨过滑轮，其一端悬挂物块 N，另一端与斜面上的物块 M 相连，系统处于静止状态。现用水平向左的拉力缓慢拉动 N，直至悬挂 N 的细绳与竖直方向成 $45^{\circ}$。已知 M 始终保持静止，则在此过程中\xzanswer{BD}

\onepicture[w=6.0cm, align=right]{\includesvg{figs/19}}

\fourchoices[answer=BD]
{水平拉力的大小可能保持不变}
{M 所受细绳的拉力大小一定一直增加}
{M 所受斜面的摩擦力大小一定一直增加}
{M 所受斜面的摩擦力大小可能先减小后增加}

%% answer: BD
%% memo:
\memoanswer{
当物块 N 受到水平拉力 $F$ 后，其受到重力、水平拉力 $F$、细绳的拉力 $T$，在物块 N 缓慢被拉动时，处于动态平衡状态。设细绳与竖直方向夹角为 $\alpha$，物块 N 的质量为 $m_1$，有 $F=T\sin\alpha$，$m_1g=T\cos\alpha$，解得 $T=\frac{m_1g}{\cos\alpha}$，$F=m_1g\tan\alpha$。在 $\alpha$ 从 $0^{\circ}$ 增大到 $45^{\circ}$ 过程中，水平拉力 $F$ 增大，细绳拉力 $T$ 增大，故 A 错误 B 正确；设物块 M 的质量为 $m_2$，斜面倾角为 $\theta$，未施加水平拉力前，若 $m_2g\sin\theta>m_1g$，则物块 M 受到的静摩擦力沿斜面向上，有 $T+f=m_2g\sin\theta$，当 $T$ 增大时，静摩擦力 $f$ 先沿斜面减小到零，然后沿斜面向下增大；若 $m_2g\sin\theta<m_1g$，物块 M 受到的静摩擦力 $f$ 沿斜面向下，有 $T=f+m_2g\sin\theta$，当 $T$ 增大时，$f$ 沿斜面向下增大，故 C 错误 D 正确。
}

\item
%% number: 20
%% paperName: 2019年全国理综Ⅰ第 20 题 6 分
%% typeId: 2
%% type: 多选题
%% chapter: 电磁感应
%% point: 法拉第电磁感应定律
%% method: 无
%% score: 6
%% degree: 450
%% duplicateId: 0
%% body:
空间存在一方向与直面垂直、大小随时间变化的匀强磁场，其边界如图（a）中虚线 MN 所示，一硬质细导线的电阻率为 $\rho$、横截面积为 $S$，将该导线做成半径为 $r$ 的圆环固定在纸面内，圆心 O 在 MN 上。$t=0$ 时磁感应强度的方向如图（a）所示，磁感应强度 $B$ 随时间 $t$ 的变化关系如图（b）所示，则在 $t=0$ 到 $t=t_1$ 的时间间隔内\xzanswer{BC}

\twopicture[w=5.5cm, align=right]{figs/20a.png}{figs/20b.png}

\fourchoices[answer=BC]
{圆环所受安培力的方向始终不变}
{圆环中的感应电流始终沿顺时针方向}
{圆环中的感应电流大小为 $\frac{B_0rS}{4t_0\rho}$}
{圆环中的感应电动势大小为 $\frac{B_0\pi r^2}{4t_0}$}

%% answer: BC
%% memo:
\memoanswer{
由楞次定律可知圆环中感应电流方向始终沿顺时针方向，故 B 正确；圆环中的磁场方向发生变化，而感应电流方向不变，根据左手定则可知圆环受到的安培力方向发生变化，故 A 错误；由法拉第电磁感应定律得 $E=\frac{\Delta\phi}{\Delta t}=\frac{\Delta B}{\Delta t}\times\frac{\pi r^2}{2}=\frac{B_0\pi r^2}{2t_0}$，故 D 错误；圆环的总电阻为 $R=\rho\frac{2\pi r}{S}$，则圆环中的感应电流 $I=\frac{E}{R}=\frac{B_0rS}{4t_0\rho}$，故 C 正确。
}

\item
%% number: 21
%% paperName: 2019年全国理综Ⅰ第 21 题 6 分
%% typeId: 2
%% type: 多选题
%% chapter: 运动和力
%% point: 变加速直线运动
%% method: 无
%% score: 6
%% degree: 350
%% duplicateId: 0
%% body:
在星球 M 上将一轻弹簧竖直固定在水平桌面上，把物体 P 轻放在弹簧上端，P 由静止向下运动，物体的加速度 $a$ 与弹簧的压缩量 $x$ 间的关系如图中实线所示。在另一星球 N 上用完全相同的弹簧，改用物体 Q 完成同样的过程，其 $a-x$ 关系如图中虚线所示，假设两星球均为质量均匀分布的球体。已知星球 M 的半径是星球 N 的 3 倍，则\xzanswer{AC}

\onepicture[w=5.5cm, align=right]{\includesvg{figs/21}}

\fourchoices[answer=AC]
{M 与 N 的密度相等}
{Q 的质量是 P 的 3 倍}
{Q 下落过程中的最大动能是 P 的 4 倍}
{Q 下落过程中弹簧的最大压缩量是 P 的 4 倍}

%% answer: AC
%% memo:
\memoanswer{
物体刚放到轻弹簧上时的加速度为星球表面的重力加速度，则 M 星球表面的重力加速度 $g_1=3a_0$，N 星球表面的重力加速度 $g_2=a_0$，在星球表面有 $G\frac{Mm}{R^2}=mg$，星球的质量为 $M=\rho\times\frac{4}{3}\pi R^3$，得 $\rho=\frac{3g}{4\pi GR}$。设 M 星球的密度为 $\rho_1$，N 星球的密度为 $\rho_2$，则有 $\frac{\rho_1}{\rho_2}=\frac{g_1}{g_2}\times\frac{R_2}{R_1}=1$，故 A 正确；设 P 物体的质量为 $m_1$，Q 物体的质量为 $m_2$，弹簧的劲度系数为 $k$，物体在弹簧上下落至加速度为零时，对 P 物体有 $kx_0=m_1g_1$，对 Q 物体有 $2kx_0=m_2g_2$，则 $\frac{m_1}{m_2}=\frac{1}{6}$，故 B 错误；将 $a-x$ 图像的纵坐标分别乘两物体的质量，则纵坐标转换为 $ma$，即物体受到的合力，图像转换为 $F_{\text{合}}-x$ 图像，图像与 $x$ 轴所围面积大小表示合力做的功的大小，由动能定理得 $W_{\text{合}}=\Delta E_k$。当加速度为零时，物体速度最大，动能最大，有 $E_{\text{km1}}=\frac{1}{2}\times 3m_1a_0\times x_0$，$E_{\text{km2}}=\frac{1}{2}\times m_2a_0\times 2x_0$，则 $\frac{E_{\text{km1}}}{E_{\text{km2}}}=\frac{3m_1}{2m_2}=\frac{1}{4}$，故 C 正确；将两图线继续延伸，当物体速度为零时合力做功为零，即 $a-x$ 图像与 $x$ 轴所围面积代数和为零，故物块 P 下降的最大距离为 $2x_0$，物块 Q 下降的最大距离为 $4x_0$，两物体下降的最大距离之比为 1:2，弹簧的最大压缩量之比为 1:2，故 D 错误。
}

\item
%% number: 22
%% paperName: 2019年全国理综Ⅰ第 22 题 5 分
%% typeId: 4
%% type: 实验题
%% chapter: 直线运动
%% point: 打点计时器公式
%% method: 无
%% score: 5
%% degree: 650
%% duplicateId: 0
%% body:
某小组利用打点计时器对物块沿倾斜的长木板加速下滑时的运动进行研究。物块拖动纸带下滑，打出的纸带一部分如图所示。已知打点计时器所用交流电的频率为 $50\UHz$，纸带上标出的每两个相邻点之间还有 4 个打出的点未画出。在 A、B、C、D、E 五个点中，打点计时器最先打出的是\tkanswer[1.2cm]{A}点，在打出 C 点时物块的速度大小为\tkanswer[2cm]{0.233} $\Ums$（保留 3 位有效数字）；物块下滑的加速度大小为\tkanswer[2cm]{0.75} $\Umsq$（保留 2 位有效数字）。

\onepicture[w=11cm, align=right]{figs/22.png}

%% answer:
\jdanswer{
\begin{enumerate}
	\item
	A

	\item
	0.233

	\item
	0.75

\end{enumerate}
}
%% memo:
\memoanswer{
物块沿长木板加速下滑，相等时间内的位移逐渐增大，故打点计时器最先打的是 A 点。打 C 点时物块的速度 $v_C=\frac{x_{BD}}{2T}=0.233\Ums$。设物块运动的加速度大小为 $a$，则有 $x_{CE}-x_{AC}=a(2T)^2$，解得 $a=0.75\Umsq$。
}

\item
%% number: 23
%% paperName: 2019年全国理综Ⅰ第 23 题 10 分
%% typeId: 4
%% type: 实验题
%% chapter: 电路
%% point: 电表改装
%% method: 无
%% score: 10
%% degree: 350
%% duplicateId: 0
%% body:
某同学要将一量程为 $250\UuA$ 的微安表改装为量程为 $20\UmA$ 的电流表。该同学测得微安表内阻为 $1200\UO$，经计算后将一阻值为 $R$ 的电阻与微安表连接，进行改装。然后利用一标准毫安表，根据图（a）所示电路对改装后的电表进行检测（虚线框内是改装后的电表）。

\twopicture[w=5.8cm, align=right]{figs/23a.png}{figs/23b.png}

\begin{enumerate}
	\item
	根据图（a）和题给条件，将（b）中的实物连接。
	\item
	当标准毫安表的示数为 $16.0\UmA$ 时，微安表的指针位置如图（c）所示，由此可以推测出改装的电表量程不是预期值，而是\tkanswer[1cm]{C}。（填正确答案标号）

	\onepicture[w=5.5cm, align=right]{figs/23c.png}

	\fourchoices[answer=C]
	{$18\UmA$}
	{$21\UmA$}
	{$25\UmA$}
	{$28\UmA$}
	\item
	产生上述问题的原因可能是\tkanswer[1cm]{AC}。（填正确答案标号）

	\fourchoices[answer=AC]
	{微安表内阻测量错误，实际内阻大于 $1200\UO$}
	{微安表内阻测量错误，实际内阻小于 $1200\UO$}
	{$R$ 值计算错误，接入的电阻偏小}
	{$R$ 值计算错误，接入的电阻偏大}
	\item
	要达到预期目的，无论测得的内阻值是否正确，都不必重新测量，只需要将阻值为 $R$ 的电阻换为一个阻值为 $kR$ 的电阻即可，其中 $k=$\tkanswer[1.5cm]{$\frac{99}{79}$}。
\end{enumerate}

%% answer:
\jdanswer{
\begin{enumerate}
	\item
	见解析

	\item
	C

	\item
	AC

	\item
	$\frac{99}{79}$

\end{enumerate}
}
%% memo:
\memoanswer{
（1）本题为对微安表的改装和校准，由题图（a）可知，缺失部分为需要改装的电路，小量程电流表改装为大量程电流表需要并联定值电阻，由此可得到实物连线如图所示。

\onepicture[w=5.5cm]{figs/23_an.png}

（2）由题图（c）可知，微安表的示数为 $160\UuA$，标准毫安表示数与微安表示数之比为 $n=\frac{16\UmA}{160\UuA}=100$。因此微安表改装后的量程为 $I=nI_g=25\UmA$。

（3）微安表满偏电流为 $I_g$，内阻为 $R_g$，改装后量程为 $I$，需并联的电阻为 $R$，则并联电阻 $R=\frac{I_gR_g}{I-I_g}$，解得 $I=I_g(1+\frac{R_g}{R})$。实际改装量程偏大，则可能是实际的微安表内阻大于 $1200\UO$，或并联电阻 $R$ 偏小，故 AC 正确。

（4）改装后的电流表实际最大测量值为 $I_1=25\UmA$ 时，并联电阻 $R=\frac{R_gI_g}{I_1-I_g}$；理论最大测量值 $I_2=20\UmA$ 时，并联电阻 $kR=\frac{R_gI_g}{I_2-I_g}$。两式联立得 $k=1.25$ 或 $k=\frac{99}{79}$。
}

\item
%% number: 24
%% paperName: 2019年全国理综Ⅰ第 24 题 12 分
%% typeId: 6
%% type: 计算题
%% chapter: 磁场
%% point: 带电粒子在磁场中的运动
%% method: 无
%% score: 12
%% degree: 450
%% duplicateId: 0
%% body:
如图，在直角三角形 OPN 区域内存在匀强磁场，磁感应强度大小为 $B$、方向垂直于纸面向外。一带正电的粒子从静止开始经电压 $U$ 加速后，沿平行于 $x$ 轴的方向射入磁场；一段时间后，该粒子在 OP 边上某点以垂直于 $x$ 轴的方向射出。已知 O 点为坐标原点，N 点在 $y$ 轴上，OP 与 $x$ 轴的夹角为 $30^{\circ}$，粒子进入磁场的入射点与离开磁场的出射点之间的距离为 $d$，不计重力。求：
\begin{enumerate}
	\item
	带电粒子的比荷；
	\item
	带电粒子从射入磁场到运动至 $x$ 轴的时间。
\end{enumerate}

\onepicture[w=6.5cm, align=right]{figs/24.png}

%% answer:
\jdanswer{
\begin{enumerate}
	\item
	$\frac{4U}{d^2B^2}$

	\item
	$\left(\frac{\pi}{2}+\frac{\sqrt{3}}{3}\right)\frac{d^2B}{4U}$ 或 $\left(\frac{\pi}{8}+\frac{\sqrt{3}}{12}\right)\frac{d^2B}{U}$

\end{enumerate}
}
%% memo:
\memoanswer{
（1）设带电粒子的质量为 $m$，电荷量为 $q$，加速后的速度大小为 $v$。由动能定理得 $qU=\frac{1}{2}mv^2$①；设粒子在磁场中做匀速圆周运动的半径为 $r$，轨迹如图所示，由洛伦兹力公式和牛顿第二定律得 $qvB=m\frac{v^2}{r}$②。

\onepicture[w=4.5cm]{figs/24_an.png}

由几何关系得 $d=\sqrt{2}r$③，联立①②③式得 $\frac{q}{m}=\frac{4U}{d^2B^2}$④。

（2）由几何关系知，带电粒子射入磁场后运动到 $x$ 轴经过的路程为 $s=\frac{\pi r}{2}+r\tan30^{\circ}$⑤，带电粒子从射入磁场到运动到 $x$ 轴的时间为 $t=\frac{s}{v}$⑥，联立②④⑤⑥式得 $t=\left(\frac{\pi}{2}+\frac{\sqrt{3}}{3}\right)\frac{d^2B}{4U}$。
}

\item
%% number: 25
%% paperName: 2019年全国理综Ⅰ第 25 题 20 分
%% typeId: 6
%% type: 计算题
%% chapter: 功和能
%% point: 动能定理
%% method: 无
%% score: 20
%% degree: 350
%% duplicateId: 0
%% body:
竖直面内一倾斜轨道与一足够长的水平轨道通过一小段光滑圆弧平滑连接，小物块 B 静止于水平轨道的最左端，如图（a）所示。$t=0$ 时刻，小物块 A 在倾斜轨道上从静止开始下滑，一段时间后与 B 发生弹性碰撞（碰撞时间极短）；当 A 返回到倾斜轨道上的 P 点（图中未标出）时，速度减为 0，此时对其施加一外力，使其在倾斜轨道上保持静止。物块 A 运动的 $v-t$ 图像如图（b）所示，图中的 $v_1$ 和 $t_1$ 均为未知量。已知 A 的质量为 $m$，初始时 A 与 B 的高度差为 $H$，重力加速度大小为 $g$，不计空气阻力。
\begin{enumerate}
	\item
	求物块 B 的质量；
	\item
	在图（b）所描述的整个运动过程中，求物块 A 克服摩擦力所做的功；
	\item
	已知两物块与轨道间的动摩擦因数均相等，在物块 B 停止运动后，改变物块与轨道间的动摩擦因数，然后将 A 从 P 点释放，一段时间后 A 刚好能与 B 再次碰上。求改变前后动摩擦因数的比值。
\end{enumerate}

\twopicture[w=6.0cm, align=right]{figs/25a.png}{figs/25b.png}

%% answer:
\jdanswer{
\begin{enumerate}
	\item
	$3m$

	\item
	$\frac{2}{15}mgH$

	\item
	$\frac{11}{9}$

\end{enumerate}
}
%% memo:
\memoanswer{
（1）根据图（b），$v_1$ 为物块 A 在碰撞前瞬间的速度大小，$\frac{v_1}{2}$ 为其碰撞后瞬间的速度大小。设物块 B 的质量为 $m'$，碰撞后的速度大小为 $v'$，由动量守恒和机械能守恒得 $mv_1=m\left(-\frac{v_1}{2}\right)+m'v'$①、$\frac{1}{2}mv_1^2=\frac{1}{2}m\left(-\frac{v_1}{2}\right)^2+\frac{1}{2}m'v'^2$②，联立①②式得 $m'=3m$③。

（2）在图（b）所描述的运动中，设物块 A 与轨道间的滑动摩擦力大小为 $f$，下滑过程中所走过的路程为 $s_1$，返回过程中所走过的路程为 $s_2$，P 点的高度为 $h$，整个过程中克服摩擦力所做的功为 $W$，由动能定理得 $mgH-fs_1=\frac{1}{2}mv_1^2-0$④、$-\left(fs_2+mgh\right)=0-\frac{1}{2}m\left(-\frac{v_1}{2}\right)^2$⑤，由图（b）所给的 $v-t$ 图线可知 $s_1=\frac{1}{2}v_1t_1$⑥、$s_2=\frac{1}{2}\times\frac{v_1}{2}\times(1.4t_1-t_1)$⑦，由几何关系可知 $\frac{s_2}{s_1}=\frac{h}{H}$⑧，物块 A 在整个过程中克服摩擦力所做的功为 $W=fs_1+fs_2$⑨，联立④⑤⑥⑦⑧⑨式可得 $W=\frac{2}{15}mgH$⑩。

（3）设倾斜轨道的倾角为 $\theta$，物块与轨道间的动摩擦因数在改变前为 $\mu$，则有 $W=\mu mg\cos\theta\times\frac{H+h}{\sin\theta}$\ccnum{⑪}；设物块 B 在水平轨道上能够滑行的距离为 $s'$，由动能定理得 $-\mu m'gs'=0-\frac{1}{2}m'v'^2$\ccnum{⑫}；设改变后的动摩擦因数为 $\mu'$，由动能定理得 $mgh-\mu'mg\cos\theta\cdot\frac{h}{\sin\theta}-\mu'mgs'=0$\ccnum{⑬}。联立以上各式可得 $\frac{\mu'}{\mu}=\frac{11}{9}$\ccnum{⑭}。
}

\item
%% number: 33
%% paperName: 2019年全国理综Ⅰ第 33 题 10 分
%% typeId: 6
%% type: 计算题
%% chapter: 气体性质
%% point: 气体多过程
%% method: 无
%% score: 10
%% degree: 550
%% duplicateId: 0
%% body:
\begin{enumerate}
	\item
	（5 分）某容器中的空气被光滑活塞封住，容器和活塞绝热性能良好，空气可视为理想气体。初始时容器中空气的温度与外界相同，压强大于外界。现使活塞缓慢移动，直至容器中的空气压强与外界相同。此时，容器中空气的温度\tkanswer[1.2cm]{低于}（填“高于”“低于”或“等于”）外界温度，容器中空气的密度\tkanswer[1.2cm]{大于}（填“大于”“小于”或“等于”）外界空气的密度。
	\item
	（10 分）热等静压设备广泛用于材料加工中。该设备工作时，先在室温下把惰性气体用压缩机压入到一个预抽真空的炉腔中，然后炉腔升温，利用高温高气压环境对放入炉腔中的材料加工处理，改善其性能。一台热等静压设备的炉腔中某次放入固体材料后剩余的容积为 $0.13\Umc$，炉腔抽真空后，在室温下用压缩机将 10 瓶氩气压入到炉腔中。已知每瓶氩气的容积为 $3.2\times10^{-2}\Umc$，使用前瓶中气体压强为 $1.5\times10^{7}\UPa$，使用后瓶中剩余气体压强为 $2.0\times10^{6}\UPa$；室温温度为 $27\UdC$。氩气可视为理想气体。
	\begin{enumerate}
		\item
		求压入氩气后炉腔中气体在室温下的压强；
		\item
		将压入氩气后的炉腔加热到 $1227\UdC$，求此时炉腔中气体的压强。
	\end{enumerate}
\end{enumerate}

%% answer:
\jdanswer{
\begin{enumerate}
	\item
	低于，大于

	\item
	$3.2\times10^{7}\UPa$，$1.6\times10^{8}\UPa$

\end{enumerate}
}
%% memo:
\memoanswer{
（1）容器绝热，活塞移动后气体压强变小，则气体的体积增大，气体对外做功，内能减小，温度降低。取与容器中质量相等的外界空气，和容器内变化后的空气相比较，二者压强相同，容器内空气的温度低于外界温度，由理想气体状态方程可知，容器内空气的体积偏小，则密度大于外界空气的密度。

（2）（i）设初始时每瓶气体的体积为 $V_0$，压强为 $p_0$；使用后气瓶中剩余气体的压强为 $p_1$。假设体积为 $V_0$、压强为 $p_0$ 的气体压强变为 $p_1$ 时，其体积膨胀为 $V_1$，由玻意耳定律得 $p_0V_0=p_1V_1$①；被压进炉腔的气体在室温和 $p_1$ 条件下的体积为 $V_1'=V_1-V_0$②；设 10 瓶气体压入完成后炉腔中气体的压强为 $p_2$，体积为 $V_2$，由玻意耳定律得 $p_2V_2=10p_1V_1'$③。联立①②③式代入题给数据得 $p_2=3.2\times10^{7}\UPa$④。

（ii）设加热前炉腔的温度为 $T_0$，加热后炉腔温度为 $T_1$，气体压强为 $p_3$。由查理定律得 $\frac{p_3}{T_1}=\frac{p_2}{T_0}$⑤，联立④⑤式并代入题给数据得 $p_3=1.6\times10^{8}\UPa$。
}

\item
%% number: 34
%% paperName: 2019年全国理综Ⅰ第 34 题 10 分
%% typeId: 6
%% type: 计算题
%% chapter: 光学
%% point: 光的折射
%% method: 无
%% score: 10
%% degree: 550
%% duplicateId: 0
%% body:
\begin{enumerate}
	\item
	（5 分）一简谐横波沿 $x$ 轴正方向传播，在 $t=\frac{T}{2}$ 时刻，该波的波形图如图（a）所示，P、Q 是介质中的两个质点。图（b）表示介质中某质点的振动图像。下列说法正确的是\xzanswer{CDE}

	\twopicture[w=6.0cm, align=right]{figs/34a.png}{figs/34b.png}

	\fivechoices[answer=CDE]
	{质点 Q 的振动图像与图（b）相同}
	{在 $t=0$ 时刻，质点 P 的速率比质点 Q 的大}
	{在 $t=0$ 时刻，质点 P 的加速度的大小比质点 Q 的大}
	{平衡位置在坐标原点的质点的振动图像如图（b）所示}
	{在 $t=0$ 时刻，质点 P 与其平衡位置的距离比质点 Q 的大}
	\item
	（10 分）如图，一艘帆船静止在湖面上，帆船的竖直桅杆顶端高出水面 $3\Um$。距水面 $4\Um$ 的湖底 P 点发出的激光束，从水面出射后恰好照射到桅杆顶端，该出射光束与竖直方向的夹角为 $53^{\circ}$（取 $\sin53^{\circ}=0.8$）。已知水的折射率为 $\frac{4}{3}$。
	\begin{enumerate}
		\item
		求桅杆到 P 点的水平距离；
		\item
		船向左行驶一段距离后停止，调整由 P 点发出的激光束方向，当其与竖直方向夹角为 $45^{\circ}$ 时，从水面射出后仍然照射在桅杆顶端，求船行驶的距离。
	\end{enumerate}

	\onepicture[w=6.5cm, align=right]{figs/34c.png}
\end{enumerate}

%% answer:
\jdanswer{
\begin{enumerate}
	\item
	CDE

	\item
	$7\Um$，$5.5\Um$

\end{enumerate}
}
%% memo:
\memoanswer{
（1）由于简谐横波沿 $x$ 轴正向传播，结合图（a）以及同侧法可知，$t=\frac{T}{2}$ 时刻 Q 质点沿 $y$ 轴正向振动，由图（b）振动图像可知，$t=\frac{T}{2}$ 时刻该质点沿 $y$ 轴负向振动，故 A 错误；由图（a）可知在 $t=\frac{T}{2}$ 时 P 在波峰位置，结合波的传播方向可知在 $t=0$ 时刻质点 P 位于波谷，此时偏离平衡位置的位移仍然最大，速率为零，加速度最大，质点 Q 仍位于平衡位置，偏离平衡位置位移为零，其振动方向与当前振动方向相反，即沿 $y$ 轴负向振动，速率最大，加速度为零，故 B 错误 CE 正确；由图（a）可知，$t=\frac{T}{2}$ 时刻平衡位置在坐标原点的质点沿 $y$ 轴负向振动，与图（b）振动情况相同，故 D 正确。

（2）（i）设光束从水面射出的点到桅杆的水平距离为 $x_1$，到 P 点的水平距离为 $x_2$；桅杆高度为 $h_1$，P 点处水深为 $h_2$；激光束在水中与竖直方向的夹角为 $\theta$，由几何关系得 $\frac{x_1}{h_1}=\tan53^{\circ}$①、$\frac{x_2}{h_2}=\tan\theta$②，由折射定律得 $\sin53^{\circ}=n\sin\theta$③；设桅杆到 P 点的水平距离为 $x$，则 $x=x_1+x_2$④，联立①②③④式并代入题给数据得 $x=7\Um$⑤。

（ii）设激光束在水中与竖直方向的夹角为 $45^{\circ}$ 时，从水面出射的方向与竖直方向的夹角为 $i'$，由折射定律得 $\sin i'=n\sin45^{\circ}$⑥；设船向左行驶的距离为 $x'$，此时光束从水面射出的点到桅杆的水平距离为 $x_1'$，到 P 点的水平距离为 $x_2'$，则 $x_1'+x_2'=x'+x$⑦、$\frac{x_1'}{h_1}=\tan i'$⑧、$\frac{x_2'}{h_2}=\tan45^{\circ}$⑨，由⑤⑥⑦⑧⑨式代入题给数据得 $x'=(6\sqrt{2}-3)\Um\approx5.5\Um$⑩。
}

\end{enumerate}

\end{document}
